\noindent\textbf{Por qué se usa este procedimiento.} Liberar el apoyo redundante y exigir desplazamiento nulo en su posición. En el voladizo base, la carga uniforme da $qL^4/(8EI)$ y la reacción ascendente $X$ da $XL^3/(3EI)$. Igualarlas permite despejar $X=3qL/8$.

\begin{enumerate}

\item La carga total es $24$ kN.

\end{enumerate}

\noindent\textbf{Operaciones y resultados intermedios.}

\begin{enumerate}

\item Reacción en el apoyo.

\[\fitmath{y_{1}=\frac{3\cdot 8\cdot 3000}{8}\approx 9000\,\mathrm{N}}\]

\end{enumerate}
